
\documentclass[conference]{IEEEtran}

\title{Labwork 1}
\author{
	Le-Chi Thanh Lam\\
	\texttt{23BI14248}
}

\begin{document}

\maketitle

\section{Introduction}

This report is an attempt at recreating the result obtained from ECG heartbeat
classficiation paper (Mohammaed et al 2018).
% Put this in the reference https://arxiv.org/abs/1805.00794

The paper discuss the classification of heartbeats to classify 5 different
arrhythmias (irregularities in heartbeat) with the AAMI EC57 standard (I don't
know what this is.) In addition, this paper also discuss a method to knowledge
transfer to myocardial infarction. I don't know how significant that is, so I
will not use the PTB Diagnostic dataset in this report.

The data is obtained from the author's Kaggle dataset.
% Link this https://www.kaggle.com/datasets/shayanfazeli/heartbeat

I limit the scope of this report to the classification of heartbeats without
discussing knowledge transfer so I will use the processed MIT BIH Arrhythmia
dataset only.

\section{Data Description}

So what's in the MIT BIH Arrhythmia dataset?

First of all this is the processed data from the original dataset.
% Reference this: https://physionet.org/content/mitdb/1.0.0/



\section{Model Implementation}

The model used in this report has a basic convolutional architecture.
The model is composed of three conv1d layers. Each have 32 kernels of size 5
with ReLU activation.
Following the convolutional layers are a max pooling layer of size 5, strides 2.
Then a flatten layer. Two fully connected layers with relu, then softmax, respectively.

\section{Training}

For training, we use the categorical cross entropy as loss function, Adam optimizer, and metrics is accuracy. 

Then we train with batch size 64, and 12 epochs.

\section{Result}

The result is 9.9785 on test loss, and 0.8276 on test accuracy.

\end{document}
